
% language=us runpath=texruns:manuals/ontarget

\startcomponent ontarget-bitmaps

\environment ontarget-style

\startchapter[title={Unusual bitmaps}]

\startsection[title=Introduction]

In the early days of \TEX, fonts mostly were bitmaps and when such a bitmap was
shown in zeros and ones, the shape was rather recognizable. A recent example of a
special|-|purpose (TAOCP) bitmap font is Don Knuths three|-|six font. Do you
recognize this character?

\starttyping
1111111111
1100110011
1100110011
0000110000
0000110000
0000110000
0001111000
0001111000
\stoptyping

It has a rather low resolution, but can still serve its purpose as we will see
later on. One problem is of course that such a low resolution doesn't render too
well. Although vector images are often preferred, especially in a \METAPOST\
context, below we will explain how we can also use bitmaps in a constructive way.

This article appeared in the first 2024 issue of \TUGBOAT\ and as usual Karl
Berry helped to make it better than it was submitted.

\stopsection

\startsection[title=Vectorizing bitmaps]

The potrace library by Peter Selinger is a nice tool: you feed it a bitmap and
are rewarded with an outline specification. Details about the process can be
found in \hyphenatedurl {https://potrace.sourceforge.net/potrace.pdf}. When you
limit yourself to only the basics, there are not that many source files and
therefore I decided to add it to \LUAMETATEX\ in order to explore if we can do
runtime conversions. Possible applications are logos and maybe converted bitmap
fonts, although these can best be prepared beforehand because consistent metrics
need to be taken care of. There are, however, other applications possible. Here I
will discuss usage only in the perspective of \METAFUN, simply because we have to
be visual, and \METAFUN\ is all about that. The library is of course accessible
from the \LUA\ end, if only because that way we can use it in \METAFUN.

We start with a simple example that also shows a potential usage:

\startbuffer
\startMPcode
    string s ; s := "010 111 010";

    draw lmt_potraced [
        bytes = s,
    ] ysized 2cm
      withpen pencircle scaled 4
      withcolor darkgreen ;
\stopMPcode
\stopbuffer

\typebuffer

We feed the \type {lmt_potraced} macro a $3\times3$ bitmap encoded as
a string and this is what we get back:

\startlocallinecorrection
\getbuffer
\stoplocallinecorrection

We expect a cross and get back a circle which is not what we want. This is
because we don't have much body in this bitmap and potrace needs more pixels in
order to give back a decent outline.

\startbuffer
\startMPcode
    string s ; s := "010 111 010";

    draw lmt_potraced [
        bytes   = s,
        explode = true,
    ] ysized 2cm
      withpen pencircle scaled 4
      withcolor darkgreen ;
\stopMPcode
\stopbuffer

\typebuffer

So, this time we \quote {explode} the bitmap, that is: we repeat every pixel
three times in the horizontal and vertical direction. One can specify \type {nx}
and \type {ny} but so far using different values doesn't help more than the magic
threesome.

\startlocallinecorrection
\getbuffer
\stoplocallinecorrection

We're getting there but need to do a bit more.

\startbuffer
\startMPcode
    string s ; s := "010 111 010";

    draw lmt_potraced [
        bytes     = s,
        threshold = 0.25,
        explode   = true,
    ] ysized 2cm
      withpen pencircle scaled 4
      withcolor darkgreen ;
\stopMPcode
\stopbuffer

\typebuffer

Here we've set a threshold which will clip the paths in a range so that our case
will get a better fit:

\startlocallinecorrection
\getbuffer
\stoplocallinecorrection

By now you will have noticed that we get back a path and this is an important
feature of potrace: it returns a closed path expressed in lines and curves that
one is supposed to fill. That result is converted into a \LUA\ table that we then
can use for instance to generate a valid \METAPOST\ path. We can do whatever we
like with that path: draw or fill it for clipping. Natively, the library might
return multiple paths but the user sees only one because we concatenate them
which is a feature of the \METAPOST\ library that comes with \LUAMETATEX.

Once we could do this it was no big deal to add support for filtering. After all,
we have more than \type {0} and \type {1} characters available. Take this
example, where we lay out the bitmap a bit differently, for clarity:

\startbuffer
\startMPcode
    string s ; s := "
        211222122
        133111311
        211222122
    ";

    path p[] ;

    p[1] := lmt_potraced [
        bytes     = s,
        threshold = 0.25,
        explode   = true,
        value     = "1",
    ] ;
    p[2] := lmt_potraced [
        bytes     = s,
        threshold = 0.25,
        explode   = true,
        value     = "2",
    ] ;
    p[3] := lmt_potraced [
        bytes     = s,
        threshold = 0.25,
        explode   = true,
        value     = "3",
    ] ;

    fill p[1] withcolor darkgreen ;
    draw p[1] withcolor darkyellow ;
    draw p[2] withcolor darkred ;
    draw p[3] withcolor darkblue ;

    currentpicture := currentpicture ysized 3cm ;
\stopMPcode
\stopbuffer

\typebuffer

We save the paths so that we can use them multiple times, here for a draw and
fill operation on the first path but you could scale, rotate, or manipulate the
result before rendering it.

\startlocallinecorrection
    \doifmodeelse {tugboat} {
        \noindent\scale[width=\textwidth]{\getbuffer}
    } {
        \noindent\getbuffer
    }
\stoplocallinecorrection

This example demonstrates that a user can define outlines using a bitmap
specification and that the amount of code is rather small. At some point we might
add a few more helpers that might reduce the amount of code even more.

\startbuffer
\startMPcode
    string s ; s := "
        212111233
        131111233
        212222133
        212222133
    ";

    path p[] ;

    p[1] := lmt_potraced [
        bytes     = s,
        threshold = 0.25,
        explode   = true,
        value     = "1",
    ] ;
    p[2] := lmt_potraced [
        bytes     = s,
        threshold = 0.25,
        explode   = true,
        value     = "2",
    ] ;
    p[3] := lmt_potraced [
        bytes     = s,
        threshold = 0.25,
        explode   = true,
        value     = "3",
    ] ;

    linejoin := butt ;

    fill p[1] withcolor darkgreen  withtransparency (1,.50) ;
    fill p[2] withcolor darkred    withtransparency (1,.50) ;
    fill p[3] withcolor darkblue   withtransparency (1,.50) ;
    draw p[1] withcolor darkyellow withtransparency (1,.75) ;
    draw p[2] withcolor darkyellow withtransparency (1,.75) ;
    draw p[3] withcolor darkyellow withtransparency (1,.75) ;

    currentpicture := currentpicture ysized 4cm ;
\stopMPcode
\stopbuffer

Because we have single paths we can safely apply properties like transparency, as
shown below: crossing lines come out right instead of with accumulated
transparent colors.

\startlocallinecorrection
    \doifmodeelse {tugboat} {
        \noindent\scale[width=\textwidth]{\getbuffer}
    } {
        \noindent\getbuffer
    }
\stoplocallinecorrection

% {\em basically we can make tabular info if we relate cells to texts}

Sometimes you want to swap the rows and columns so we provide a feature for doing
that:

\startbuffer
\startMPcode
    string s[] ;

    s[1] := "1110 0110 0110 0111";
    s[2] := "1000 1111 1111 0001";

    draw lmt_potraced [
        bytes   = s[1],
        explode = true,
    ] ysized 2cm withcolor darkgreen withpen pencircle scaled 4 ;

    draw lmt_potraced [
        bytes   = s[2],
        explode = true,
    ] ysized 2cm withcolor darkblue withpen pencircle scaled 4 ;

    draw lmt_potraced [
        bytes   = s[1],
        explode = true,
        swap    = true,
    ] ysized 2cm withcolor white withpen pencircle scaled 2 ;
\stopMPcode
\stopbuffer

\typebuffer

When one uses \LUA\ input (as we will see later), one can do that when generating
the bitmap. At any rate, this is what we get from the above:

\startlocallinecorrection
\getbuffer
\stoplocallinecorrection

The previous examples demonstrate that not much code is needed in order to
achieve nice effects. It also illustrates that one needs to twist the mind a
little and think of bitmap specifications as actually efficient outline
definitions. One could argue that such rectangular shapes are easy to program in
\METAPOST\ anyway, and going via bitmaps is kind of strange. So, let's move on to
a more attractive example.

\stopsection

\startsection[title=How about fonts]

In order to demonstrate building fonts we need a decent bitmap font and it
happens that Don Knuth's \quote {Font36} is a good candidate. We have discussed
that one elsewhere and its usage can be found in the \METAFUN\ \type {threesix}
library. We happily borrow the definitions from that library:

\startbuffer
\startMPdefinitions
string dekthreesix[] ; path shapes[] ;

def DEK(expr n, b) =
    dekthreesix[utfnum(n)] := b ;
enddef ;

DEK("0", "00111100 01111110 11000011 11000011 11000011 11000011 01111110 00111100" ) ;
DEK("1", "00011100 11111100 11101100 00001100 00001100 00001100 11111111 11111111" ) ;
DEK("2", "00111110 01110011 01110011 00000011 00001110 00111000 11110001 11111111" ) ;
DEK("3", "01111110 01100110 00000110 00111100 00000110 11100111 11100111 01111110" ) ;
DEK("4", "00001110 00011110 00110110 01100110 11111111 11111111 00000110 00001110" ) ;
DEK("5", "01111110 01111110 01000000 01111100 01000111 00000011 11000111 11111110" ) ;
DEK("6", "01111110 01100110 11000000 11011100 11100110 11000011 01100011 01111110" ) ;
DEK("7", "11111111 11111111 10000111 10001110 00011100 00011100 00011100 00011100" ) ;
DEK("8", "01111110 01100110 01100110 00111100 11000011 11000011 11000011 01111110" ) ;
DEK("9", "01111110 11000110 11000011 01100111 00111011 00000011 01100110 01111110" ) ;
DEK("A", "000110000 000111000 001111100 001101100 011001110 011111110 010000110 111100111") ;
DEK("B", "1111100 1110110 0110110 0111100 0110011 0110011 1110011 1111100") ;
DEK("C", "01111101 11110011 11100001 11100001 11100000 11100000 11100001 01111110") ;
DEK("D", "11111100 11100010 01100011 01100011 01100011 01100011 11100010 11111100") ;
DEK("E", "1111111 1110001 0110101 0111100 0110100 0110001 1110001 1111111") ;
DEK("F", "11111111 11100001 01100101 01111100 01100100 01100000 11111000 11111000") ;
DEK("G", "01111010 11100110 11100010 11100000 11100111 11100010 11100010 01111100") ;
DEK("H", "11100111 11100111 01000010 01111110 01000010 01000010 11100111 11100111") ;
DEK("I", "1111111 1111111 0011000 0011000 0011000 0011000 1111111 1111111") ;
DEK("J", "01111111 01111111 00000110 00000110 11110110 01100110 01100110 00111100") ;
DEK("K", "11101110 11100100 01101000 01110000 01111000 01101100 11100110 11101111") ;
DEK("L", "111111000 111111000 011000000 011000000 011000000 011000011 111000011 111111111") ;
DEK("M", "1100000011 1110000111 0111111110 0100110010 0100110010 0100000010 1100000011 1100000011") ;
DEK("N", "11000011 11100011 01110010 01111010 01011110 01001110 11100110 11100010") ;
DEK("O", "01111110 11000011 11000011 11000011 11000011 11000011 11000011 01111110") ;
DEK("P", "11111110 11100011 01100011 01111110 01100000 01100000 11111000 11111000") ;
DEK("Q", "01111100 11000110 11000110 11000110 11000110 11001110 11000110 01111101") ;
DEK("R", "11111100 11100110 01100110 01111100 01101000 01100100 11100010 11100111") ;
DEK("S", "01111110 11100001 11100001 01111000 00011110 10000111 11000011 10111110") ;
DEK("T", "1111111111 1100110011 1100110011 0000110000 0000110000 0000110000 0001111000 0001111000") ;
DEK("U", "111101111 111101111 011000010 011000010 011000010 011000010 011000010 001111100") ;
DEK("V", "11111000111 11111000111 01110000010 00110000100 00111000100 00011001000 00001101000 00001110000") ;
DEK("W", "111000000111 111000000111 110000000010 011000000100 011001000100 001101101000 001101101000 000110110000") ;
DEK("X", "1111001110 1111000110 0001101000 0000110000 0000110000 0001011000 0110001111 0111001111") ;
DEK("Y", "111100011 111100011 011000010 001110100 000111000 000011000 001111110 001111110") ;
DEK("Z", "11111111 10000111 00001110 00011100 00111000 01110000 11100001 11111111") ;
\stopMPdefinitions
\stopbuffer

\doifelsemode {tugboat} {
    \typebuffer
} {
    \typebuffer[style=\small\tt]
}

\getbuffer

We use these definitions in the \METAPOST\ code below. Helpers like \type
{utfnum} are part of \LUAMETAFUN\ and we use (named) colors defined at the
\CONTEXT\ end.

\startbuffer
\startMPcode
def shapethem(expr first, last) =
    for i = utfnum(first) upto utfnum(last) :
        shapes[i] := lmt_potraced [
            bytes   = dekthreesix[i],
            explode = true,
          % threshold = .25, % more rectangular
            value   = "1",
        ] ;
    endfor ;
enddef ;

def drawthem(expr first, last, dx, dy) =
    numeric d ; d := 0 ;
    numeric f ; f := utfnum(first) ;
    numeric l ; l := utfnum(last) ;
    for i = f upto l :
        fill shapes[i] shifted (d,dy) withcolor "middlegray" ;
        draw shapes[i] shifted (d,dy) withcolor "darkgreen" ;
      % draw boundingbox shapes[i] shifted (d,dy);
        d := d + bbwidth(shapes[i]) + dx;
    endfor ;
enddef ;

shapethem("0","9") ;
shapethem("A","Z") ;

drawthem("0", "9", 10, -00) ;
drawthem("A", "J", 10, -30) ;
drawthem("K", "S", 10, -60) ;
drawthem("T", "Z", 10, -90) ;

currentpicture := currentpicture xsized TextWidth ;
\stopMPcode
\stopbuffer

\typebuffer

Here we don't integrate it as a font but just show the characters as they come
out, which hopefully is easier to understand. You can take a close look at the
bitmaps in order to see what we get rendered in \in {figure} [fig:fancyfont].
Setting a lower threshold will give more rectangular results.

\startplacefigure[title=A fancy outline font,reference=fig:fancyfont]
    \getbuffer
\stopplacefigure

Because we're not going to use this font for typesetting here, a simple example
of a definition will do. We first load an already existing module so that we
don't need to define the bitmap definitions; basically they are the same as
above.

\startbuffer
\useMPlibrary[threesix]

\startMPcalculation{simplefun}
    vardef ThreeSixPotraced (expr code, spread, lift) =
        draw lmt_potraced [
            bytes   = code,
            explode = true,
            value   = "1",
        ] scaled (1/3) shifted (spread,lift) ;
    enddef ;
\stopMPcalculation
\stopbuffer

\typebuffer \getbuffer

Next we register a \METAFUN\ font using similar trickery as in that module. There
we define a few variants but that is not needed here.

\startbuffer
\startluacode
    local utfbyte = utf.byte

    local f_code = string.formatters['ThreeSixPotraced("%s",%s,%s);']

    function MP.registerthreesixpotraced(name)
        fonts.dropins.registerglyphs {
            name     = name,
            units    = 12,
            usecolor = true,
        }
        for u, data in table.sortedhash(MP.font36) do
            local ny     = 8
            local nx     = ((#data + 1) // ny) - 1
            local height = ny * 1.1 - 0.1
            local width  = nx * 1.1 - 0.1
            local spread = 0.9
            local lift   = 0.3
            fonts.dropins.registerglyph {
                category = name,
                unicode  = utfbyte(u),
                width    = width + spread,
                height   = height,
                code     = f_code(data,spread,lift),
            }
        end
    end

    MP.registerthreesixpotraced("fontthreesixpotraced")
\stopluacode
\stopbuffer

\typebuffer \getbuffer

Finally we define a font feature that will hook the previous code into the font
handler. There a \TYPETHREE\ font will be constructed.

\startbuffer
\definefontfeature
  [fontthreesixpotraced]
  [default]
  [metapost=fontthreesixpotraced,
   spacing=.375 plus .2 minus .1 extra .375]

\definefont[DEKFontP][Serif*fontthreesixpotraced]
\stopbuffer

\typebuffer \getbuffer

We now show an example of usage (abridged). This font only has uppercase
characters so one might consider duplicating these into the lowercase slots.
Because we replace glyphs in the serif font used, we still have a complete font,
albeit of mixed design.

% {\DEKFontP \showglyphs THIS IS JUST A TEST}
\startbuffer
\DEKFontP \WORD{\samplefile{knuth}}
\stopbuffer

\typebuffer

This gives us a rendering that is quite readable, especially when you consider
how small the bitmaps are (the red bar is the overfull box marker):

\startbuffer[tex]
\protected\def\TeX
  {\dontleavehmode
   \begingroup
   T%
   \kern-.40\fontcharwd\font`T%
   \lower.45\fontcharht\font`X\hbox{E}%
   \kern-.15\fontcharwd\font`X%
   X%
   \endgroup}
\stopbuffer

\start \pushoverloadmode
    \def|#1|{\endash}
    \getbuffer[tex]
    \doifelsemode {tugboat} {
        \raggedright
        \blank
        \getbuffer
        \blank
    } {
        \getbuffer
    }
\stop

Just in case Don Knuth runs into this example, we need to cheat a little here and
redefine the \TEX\ logo definition:

\typebuffer[tex]

A real font would have proper font kerns but if needed you can set that up using
the \OPENTYPE\ feature plug in mechanism that \CONTEXT\ provides. A font like
this can also be colored:

\startlocallinecorrection
\scale[height=2cm]{\darkblue\pushoverloadmode\DEKFontP\getbuffer[tex]\TeX}
\stoplocallinecorrection

\stopsection

\startsection[title=External bitmaps]

A convenient way to include traced images is to use the potrace command line
tool. However, we can also include grayscale single|-|byte \PNG|-|encoded images:

\startbuffer
\startMPcode
    path p ; p := lmt_potraced [
        filename  = "mill.png",
        criterium = 100,
    ] ;

    fill last_potraced_bounds withcolor "middlegray" ;
    fill p withcolor "darkgreen" ;
    draw p withcolor "darkred" withpen pencircle scaled 1.5 ;
    setbounds currentpicture to last_potraced_bounds ;

    currentpicture := currentpicture ysized 10cm ;
\stopMPcode
\stopbuffer

\typebuffer

This is not the best image but the photograph happens to be part of the \CONTEXT\
distribution so why not use it. You can of course apply a different criterion and
overlay a subsequent trace in a different color. The result is shown in \in {figure}
[fig:tracedmill].

\startplacefigure[title=A traced \PNG\ image.,reference=fig:tracedmill]
    \getbuffer
\stopplacefigure

In addition to the \typ {last_potraced_bounds} path variable we also have \typ
{last_potraced_width} and \typ {last_potraced_height} numeric available.

\stopsection

\startsection[title=Using \LUA|-|generated bitmaps]

It is tempting to see what can be done with a bitmap generated by \LUA, but let's
first give a simple example. Here we register a bitmap:

\startbuffer
\startluacode
local s = [[
    000000010000000000000010000000
    000000101000000000000101000000
    000001000100000000001000100000
    000010000010000000010000010000
    000100000001000000100000001000
    001000000000100001000000000100
    010000000000010010000000000010
    100000000000001100000000000001
]]

potrace.setbitmap("mybitmap",s)
\stopluacode
\stopbuffer

\typebuffer \getbuffer

\startbuffer
\startMPcode
    path p ; p := lmt_potraced [
        stringname = "mybitmap",
    ] ;

    fill p ysized 2cm withcolor "darkblue" ;
\stopMPcode
\stopbuffer

\typebuffer

We could play with the parameters but what we get is an outline that is supposed
to be filled:

\startlocallinecorrection
    \doifmodeelse {tugboat} {
        \noindent\scale[width=\textwidth]{\getbuffer}
    } {
        \noindent\getbuffer
    }
\stoplocallinecorrection

Next we create a bitmap from a dataset; in this case, we draw a function. Of
course we could create some handy helpers to do this:

\startbuffer
\startluacode
local d = table.setmetatableindex("table")

local t    = { }
local step = 100
local ymin = 0
local ymax = 0

local x    = 0
local dx   = 10*math.pi/step

for i=1,step do
    local y = math.round(math.cos(x)*20)
    x = x + dx
    if y > ymax then
        ymax = y
    end
    if y < ymin then
        ymin = y
    end
    t[i] = y
end

for i=1,step do
    for j=ymin, ymax do
        d[j][i] = '0'
    end
end

for i=1,step do
    d[t[i]][i] = '1'
end

for y=ymin,ymax do
    d[y] = table.concat(d[y])
end

potrace.setbitmap("mybitmap",table.concat(d," ",ymin,ymax))
\stopluacode
\stopbuffer

\typebuffer \getbuffer

\startbuffer
\startMPcode
    path p ; p := lmt_potraced [
        stringname = "mybitmap",
        explode   = true,
      % tolerance = 0.5,
        threshold = 0,
      % optimize  = true,
    ] ;

    fill p withcolor "darkblue" ;
\stopMPcode
\stopbuffer

% if texmode("tugboat") : currentpicture := currentpicture xsized(TextWidth) fi

\typebuffer

To what extent this is a useful application is to be decided:

\startlocallinecorrection
    \doifmodeelse {tugboat} {
        \noindent\scale[width=\textwidth]{\getbuffer}
    } {
        \noindent\getbuffer
    }
\stoplocallinecorrection

Later we will see that it is less work if we stay in the first quadrant, using
$(1,1)$ as lower left corner. Here we explicitly need to provide \type {concat}
the range because we have a negative $y_{\tf min}$.

It quickly gets more interesting when we use an intermediate bitmap for (a kind
of) surface plots.

% \startbuffer
% \startluacode
% local sin, cos, pi = math.sin, math.cos, math.pi
%
% local d  = table.setmetatableindex("table")
% local pp = pi/10
%
% for x=1,100 do
%     for y=1,100 do
%         local z = sin(pp*x) + cos(pp*y)
%         if z > 0.5 then
%             d[x][y] = '1'
%         elseif z > 0 then
%             d[x][y] = '2'
%         elseif z < -0.5 then
%             d[x][y] = '3'
%         else
%             d[x][y] = '4'
%         end
%     end
% end
%
% potrace.setbitmap("mybitmap",d)
% \stopluacode
% \stopbuffer

\startbuffer
\startluacode
local sin, cos, pi = math.sin, math.cos, math.pi

local pp = pi/10

local function f(x,y)
    local z = sin(pp*x) + cos(pp*y)
    if z > 0.5 then
        return '1'
    elseif z > 0 then
        return '2'
    elseif z < -0.5 then
        return '3'
    else
        return'4'
    end
end

potrace.setbitmap("mybitmap",potrace.contourplot(100,100,f))
\stopluacode
\stopbuffer

\typebuffer \getbuffer

Here we use a helper that runs the given function over the maxima and collects
the results in a bitmap. More such helpers will be provided when users come up
with more demands.

\startbuffer[1]
\startMPcode
    fill lmt_potraced [
        stringname = "mybitmap",
        value      = "1",
        explode    = true,
        threshold  = 0.25,
      % tolerance  = 0.1,
      % threshold  = 1.0,
        optimize   = true,
    ] withcolor "darkred" ;

    fill lmt_potraced [
        stringname = "mybitmap",
        value      = "2",
        explode    = true,
        threshold  = 0.25,
      % tolerance  = 0.1,
      % threshold  = 1.0,
        optimize   = true,
    ] withcolor "darkgreen" ;

    fill lmt_potraced [
        stringname = "mybitmap",
        value      = "3",
        explode    = true,
        threshold  = 0.25,
      % tolerance  = 0.1,
      % threshold  = 1.0,
        optimize   = true,
    ] withcolor "darkblue" ;

    fill lmt_potraced [
        stringname = "mybitmap",
        value      = "4",
        explode    = true,
        threshold  = 0.25,
      % tolerance  = 0.1,
      % threshold  = 1.0,
        optimize   = true,
    ] withcolor "darkyellow" ;

    clip currentpicture to last_potraced_bounds enlarged -1;
    currentpicture := currentpicture xysized (.45*TextWidth,4cm) ;
\stopMPcode
\stopbuffer

\startbuffer[2]
\startMPcode
    fill lmt_potraced [
        stringname = "mybitmap",
        value      = "1",
        explode    = true,
      % threshold  = 0.25,
        tolerance  = 0.1,
        threshold  = 1.0,
        optimize   = true,
    ] withcolor "darkred" ;

    fill lmt_potraced [
        stringname = "mybitmap",
        value      = "2",
        explode    = true,
      % threshold  = 0.25,
        tolerance  = 0.1,
        threshold  = 1.0,
        optimize   = true,
    ] withcolor "darkgreen" ;

    fill lmt_potraced [
        stringname = "mybitmap",
        value      = "3",
        explode    = true,
      % threshold  = 0.25,
        tolerance  = 0.1,
        threshold  = 1.0,
        optimize   = true,
    ] withcolor "darkblue" ;

    fill lmt_potraced [
        stringname = "mybitmap",
        value      = "4",
        explode    = true,
      % threshold  = 0.25,
        tolerance  = 0.1,
        threshold  = 1.0,
        optimize   = true,
    ] withcolor "darkyellow" ;

    clip currentpicture to last_potraced_bounds enlarged -1;
    currentpicture := currentpicture xysized (.45*TextWidth,4cm) ;
\stopMPcode
\stopbuffer

\typebuffer[1]

Here we are only exploring the possibilities so there is no interface like we
have with contour plots. One can imagine that we provide this as a plugin, which
is not that hard but whether it eventually happens depends on user demand or
rainy days. So to summarize: here we generate the bitmap, we call out to potrace
for a vector representation, that gets fed into \METAPOST\ and which gives us
back a result that can be converted to \PDF. Of course we could go directly from
potrace output to \PDF\ if we want to, but now we get full control over the final
result. In case you wonder about performance: it compiles real fast!

\startlocallinecorrection
    \noindent\getbuffer[1]\enspace\getbuffer[2]
\stoplocallinecorrection

Some work is needed to scale the image to the proportions that reflect the input
ranges but because we have a vector image that does not affect the quality of the
outcome. Here we also show the effects of \typ {tolerance} which influences the
optimizing and \typ {threshold} that determines the accuracy (number of points).
In the second rendering: we use the commented values that work quite well with
this kind of more mathematical images.

\startsection[title=Experimenting]

You can use bitmaps as a design tool but it needs a little experimenting to get
the idea. Take these examples:

\startbuffer[1]
\startMPcode
    fill lmt_potraced [ bytes = "
        00100000001
        01010000010
        10001000100
        01000101000
        00100010000
        00010101000
        00001000100
        00010100010
        00100010001
        01000001010
        10000000100
    " ] xsized 2cm
        withcolor darkgreen ;
\stopMPcode
\stopbuffer

\startbuffer[2]
\startMPcode
    fill lmt_potraced [ bytes = "
        00100000001
        01010000011
        10001000100
        01000101000
        00100010000
        00010101000
        00001000100
        00010100010
        00100010001
        11000001010
        10000000100
    " ] xsized 2cm
        withcolor darkblue ;
\stopMPcode
\stopbuffer

\startbuffer[3]
\startMPcode
    fill lmt_potraced [ bytes = "
        00100000011
        01110000111
        11111001110
        01111111100
        00111111000
        00011111000
        00011111100
        00111111110
        01110011111
        11100001110
        11000000100
    " ] xsized 2cm
        withcolor darkyellow ;
\stopMPcode
\stopbuffer

\startlocallinecorrection
\startcombination[nx=3,ny=1]
    {\getbuffer[1]} {}
    {\getbuffer[2]} {}
    {\getbuffer[3]} {}
\stopcombination
\stoplocallinecorrection

We started with a simple X|-|like symbol:

\typebuffer[1]

Next we started filling the shape a bit and tried to make it less spiked:

\typebuffer[2]

And finally we add even more ones to the bitmap. You need to fill a bitmap area
in order to get an efficient fill.

\typebuffer[3]

If you're in doubt you can also render the bitmap, assuming that it has
reasonable proportions.

\startbuffer[1]
\startMPcode
    string s ; s := "
        00100000011
        01110000111
        11111001110
        01111111100
        00111111000
        00011111000
        00011111100
        00111111110
        01110011111
        11100001110
        11000000100
    " ;
    fill lmt_potraced [ bytes = s ]
        xsized 3cm
        withcolor darkyellow ;
    draw lmt_potraced [ bytes = s, alternative = "text" ]
        shifted (.25,.25)
        xsized 3cm ;
\stopMPcode
\stopbuffer

\startbuffer[2]
\startMPcode
    string s ; s := "
        00100000011
        01110000111
        11111001110
        01111111100
        00111111000
        00011111000
        00011111100
        00111111110
        01110011111
        11100001110
        11000000100
    " ;
    fill lmt_potraced [ bytes = s, polygon = true ]
        shifted (.25,.25)
        xsized 3cm
        withcolor darkblue ;
    fill lmt_potraced [ bytes = s ]
        xsized 3cm
        withtransparency (1,.75)
        withcolor darkyellow ;
    draw lmt_potraced [ bytes = s, alternative = "text" ]
        shifted (.25,.25)
        xsized 3cm ;
\stopMPcode
\stopbuffer

\typebuffer[1]

The article that we mentioned in the introduction explains how potrace looks at a
bits in relation to its neighbors.

\startlocallinecorrection
    \noindent\getbuffer[1]\quad\getbuffer[2]
\stoplocallinecorrection

You can save some runtime (and coding) by using the start|-|stop wrappers that
keep the (intermediate) potrace object available. This permits for instance
showing the \quote {original} polygon that serves as basis for successive steps
in the library towards to the final curve(s).

\startbuffer
\startMPcode
    string s ; s :=
        "01111111111111111111111111111100
         11000000000000000000000000000110
         11000000000000000000000000000011
         11000000000000000000000000000011
         11000000000000000000000000000011
         01100000000000000000000000000011
         00111111111111111111111111111110";

    lmt_startpotraced [ bytes = s ] ;

        p := lmt_potraced [
            value     = "0",
            threshold = 0,
            tolerance = 0,
            optimize  = true,
        ] ;

        draw image (
            p := p shifted - center p ;
            draw p
                withpen pencircle scaled 1
                withcolor "darkblue" ;
            drawpoints p
                withpen pencircle scaled .5
                withcolor "white" ;
            draw boundingbox p
                withpen pencircle scaled .1
                withcolor "darkgray" ;
        ) ysized 3cm ;

        path p ; p := lmt_potraced [
            value   = "0",
            polygon = true,
        ] ;

        draw image (
            p := p shifted - center p ;
            draw p
                withpen pencircle scaled .25
                withcolor "middleyellow" ;
            drawpoints p
                withpen pencircle scaled .175
                withcolor "white" ;
            draw boundingbox p
                withpen pencircle scaled .05
                withcolor "darkgray" ;
        ) ysized 3cm ;

    lmt_stoppotraced ;
\stopMPcode
\stopbuffer

\typebuffer

\startlocallinecorrection
    \doifmodeelse {tugboat} {
        \noindent\scale[width=\textwidth]{\getbuffer}
    } {
        \noindent\getbuffer
    }
\stoplocallinecorrection

The above is just a bit of exploring the possibilities so eventually there will
be a chapter on this in the \LUAMETAFUN\ manual, because it is definitely fun to
play with this in the perspective of \METAPOST.

\stopsection

\startsection[title=Contour plots]

We end with showing how we can do rather nice contour plots and region plots with
help of the potracer. Let us start with the latter type of graphics. In a recent
math paper Mikael was counting nodal domains of Neumann eigenfunctions to the
Laplace operator in a square. These eigenfunctions are built from cosines. One
example is given by

\startformula
    \Psi(x,y) = \cos(8\pi x)\cos(3\pi y) + \cos(3\pi y)\cos(8\pi x)\mtp{.}
\stopformula

The related graphic of interest is to fill the part of the unit square where \im
{\Psi} is positive. In the article, Wolfram Mathematica was used to produce this
graphic, with its built-in function \typ {RegionPlot}. The result was indeed
satisfactory (\in{Figure}{(a)}[fig:Mathematica]).

If one looks closely at the graphics one will see that the filled regions are
made up of a mesh. With potrace we instead receive one~(!) path. To generate the
corresponding graphic we first use \LUA\ to define a function that takes a point
as input and returns 1 if \im {\Psi} is positive at the point and 0 otherwise. We
then use it to generate a \im {1000\times 1000} bitmap image of zeros and ones
accordingly.

\startbuffer
\startluacode
    local cos, pi = math.cos, math.pi

    local N       = 1000
    local pp      = pi/N
    local pp3     = 3 * pp
    local pp8     = 8 * pp
    local cospp8y = 0
    local cospp3y = 0

    local function f(x,y)
        if x == 1 then
            cospp8y = cos(pp8*y)
            cospp3y = cos(pp3*y)
        end
        local z = cos(pp8*x)*cospp3y + cos(pp3*x)*cospp8y
        if z > 0 then
            return '1'
        else
            return '0'
        end
    end

    potrace.setbitmap("mybitmap", potrace.contourplot(N,N,f))
\stopluacode
\stopbuffer

\typebuffer
\getbuffer

Once this is done, we can use the result in \typ{lmt_potraced}.

\startbuffer
\startMPcode
    path p ; p := lmt_potraced [
        stringname = "mybitmap",
        value      = "1",
        threshold  = 0.25,
        optimize   = true,
    ] ;

    p := p xsized .45TextWidth ;
    fill p withcolor "darkred" ;
\stopMPcode
\stopbuffer

\typebuffer

This results in \in{Figure}{(b)}[fig:Mathematica], which looks very similar to
the one generated by Wolfram Mathematica. Note that \type {p} is one
(disconnected) path.

\startplacefloat
    [figure]
    [reference=fig:Mathematica,
     title={(a) A region plot generated by Wolfram Mathematica.
            (b) The same, generated by potrace and \METAFUN.}]
    \startcombination[nx=2,ny=1,distance=.02tw]
        {\externalfigure[ontarget-bitmaps-mathematica][width=.45tw]}{(a)}
        {\getbuffer}{(b)}
    \stopcombination
\stopplacefloat

We give one example of a contour plot. By a contour plot, here we mean a plot of
the curve that is described as the solution to an equation \im {F(x,y) = 0}. With
\im {F(x,y) = y - f(x)} we realize that function graphs \im {y = f(x)} provide a
particular example, but we can also handle more complicated curves here, for
example the unit circle (\im {F(x,y) = x^2 + y^2 -1}).

Let us draw a part of the curve that goes under the name the Trisectrix of
Maclaurin, a curve that can be used to trisect angles (named after Colin
Maclaurin in 1742). We use the function

\startformula
    F(x,y) = 2x(x^2 + y^2) - (3x^2 - y^2)\mtp{.}
\stopformula

Let us construct the path and draw it.

\startbuffer
\startluacode
    local N   = 1000
    local xx  = 3/N
    local yy  = 3/N

    local function f(x,y)
        local x = xx*x - 1
        local y = yy*y - 1.5
        local z = 2*x*(x^2 + y^2) - (3*x^2 - y^2)
        if z > 0 then
            return '1'
        else
            return '0'
        end
    end

    potrace.setbitmap("mybitmap", potrace.contourplot(N,N,f))
\stopluacode
\stopbuffer

\typebuffer

\getbuffer

\startbuffer[1]
\startMPcode
    path p ; p := lmt_potraced [
        stringname = "mybitmap",
        value      = "1",
        tolerance  = 0.1,
        threshold  = 1,
        optimize   = true,
    ] ;
    p := p xsized TextWidth ;
    draw p withcolor "darkred" ;
    drawpoints p withcolor "orange" ;
    drawpointlabels p ;
    currentpicture := currentpicture ysized .4TextHeight ;
\stopMPcode
\stopbuffer

\typebuffer[1]

The leftmost graphic in \in {figure} [fig:contourplot] shows the result of this
code. We also draw the points and the point labels. This way we can easily find
out which part of the path to draw. In this case it seems that we need the first
\im {34} points.

\startbuffer[2]
\startMPcode
    path p ; p := lmt_potraced [
        stringname = "mybitmap",
        value      = "1",
        tolerance  = 0.1,
        threshold  = 1,
        optimize   = true,
    ] ;
    p := subpath(0,33) of p ;
    p := p xsized 4cm ;
    draw p
        withcolor "darkred"
        withpen pencircle scaled .5mm ;
    currentpicture := currentpicture ysized .4TextHeight ;
\stopMPcode
\stopbuffer

\typebuffer[2]

You need to keep the resolution (determined by the input) and therefore scaling
in mind and choose the pen accordingly as in the second drawing in
\in {figure} [fig:contourplot].

\startplacefigure[title=The same contour plot with different details.,reference=fig:contourplot]
    \startcombination[nx=2,ny=1,distance=.05tw]
        {\getbuffer[1]} {}
        {\getbuffer[2]} {}
    \stopcombination
\stopplacefigure

This method of drawing contour plots is efficient, and it gives, in contrast to
some traditional methods, curves that look smooth, with relatively few points. It
also has weaknesses, one of them being that the inequality \im {F(x,y) > 0} does
not always single out the curve \im {F(x,y) = 0}; it might be the case that the
function \im {F} is positive on both sides of the curve \im {F(x,y) = 0}.

We invite the reader to be creative and play with this (to us) new toy. Have fun!

\stopsection

\startsection[title=Coda]

And, in the spirit of having fun, here are some final images. This first one is a
photo of a \quote {tool} that we came up with at the \CONTEXT\ meeting. Willi
Egger made a kit for the attendees, so there was the usual cutting and glueing
involved. This kit is shown in \in {figure} [fig:bitmapdevice]. The dots are
seeds. It also shows an emoji that Mikael's children made when we were playing
with this feature. The first image has the rendering of that input, while the
second \quote {explodes} the pixels in the x direction (so columns are
duplicated), resulting in a more symmetric image.

\startbuffer[1]
draw image (
    lmt_startpotraced [
        bytes =
            "..11111..
             .1.....1.
             1.......1
             1.22.22.1
             1.......1
             1.3...3.1
             1..333..1
             .1.....1.
             ..11111.."
    ] ;
    fill lmt_potraced
        [ value = "1", size = 1 ]
        withcolor "darkgray"
        withpen pencircle scaled 0.1 ;
    fill lmt_potraced
        [ value = "2", size = 1 ]
        withcolor "darkblue"
        withpen pencircle scaled 0.1 ;
    fill lmt_potraced
        [ value = "3", size = 1 ]
        withcolor "darkred"
        withpen pencircle scaled 0.1 ;
    lmt_stoppotraced ;
) sized 5cm ;
\stopbuffer

\startbuffer[2]
draw image (
    lmt_startpotraced [
        nx    = 2, % more symmetric, 3 gives different curves
     %  ny    = 3, % clown
        bytes =
            "..11111..
             .1.....1.
             1.......1
             1.22.22.1
             1.......1
             1.3...3.1
             1..333..1
             .1.....1.
             ..11111.."
    ] ;
    fill lmt_potraced
        [ value = "1", size = 1 ]
        withcolor "darkgray"
        withpen pencircle scaled 0.1 ;
    fill lmt_potraced
        [ value = "2", size = 1 ]
        withcolor "darkblue"
        withpen pencircle scaled 0.1 ;
    fill lmt_potraced
        [ value = "3", size = 1 ]
        withcolor "darkred"
        withpen pencircle scaled 0.1 ;
    lmt_stoppotraced ;
) sized 5cm ;
\stopbuffer

\startplacefigure[title=A bitmap kit by Willi Egger.,reference=fig:bitmapdevice]
    \startcombination[nx=3,ny=1,distance=5mm]
        {\externalfigure[ontarget-bitmaps-device][height=5cm]} {}
        {\processMPbuffer[1]} {}
        {\processMPbuffer[2]} {}
    \stopcombination
\stopplacefigure

We only show the second emoji. When we set \type {ny} too, we get a still
interesting but somewhat weird face instead. You can try that yourself.

\typebuffer[2]

\stopsection

\stopchapter

% % >mutool convert -O colorspace=gray,resolution=200 -F png temp-potracex.pdf
%
% % \startMPpage[offset=1bp]
% %     draw for i=0 step 2pi/50 until 2pi :
% %         (2*i, sin(i)^2+sqrt(i)) ..
% %     endfor nocycle withpen pencircle scaled 0 ;
% % \stopMPpage
%
% \startMPpage
%
%     path p ; p := lmt_potraced [
%         filename  = "out-0001.png",
%         criterium = 255,
%         optimize = true,
% %         explode   = true,
% %         tolerance = 0,
% %         threshold = 0,
%     ] ;
%
%     fill p withcolor "darkred" ;
%
% \stopMPpage

\stopcomponent

\startluacode
local sin, cos, pi = math.sin, math.cos, math.pi

local N   = 1000
local pp  = pi/N
local yy  = 1/N

local function f(x,y)
    local z = yy*y - cos(pp*x)
    if z > 0 then
        return '1'
    else
        return '0'
    end
end

potrace.setbitmap("mybitmap", potrace.contourplot(N,N,f))
\stopluacode

\startMPpage[offset=1TS]
    path p ; p := lmt_potraced [
        stringname = "mybitmap",
        value      = "1",
        threshold  = 1,
        tolerance  = 0.01,
        optimize   = true,
    ] ;

    draw p withpen pencircle scaled 3 withcolor "darkred" ;

    currentpicture := currentpicture xysized (4cm,4cm) ;
    path cp ; cp := boundingbox currentpicture ;

    drawarrow ulcorner cp -- llcorner cp ;
    drawarrow ulcorner cp -- urcorner cp ;

    label.urt ("\im{x}",llcorner cp) ;
    label.llft("\im{y}",urcorner cp) ;
\stopMPpage

% \startluacode
% local sin, cos, pi = math.sin, math.cos, math.pi
%
% local N       = 1000
% local pp      =     pi/N
% local pp3     = 3 * pp
% local pp8     = 8 * pp
% local cospp8y = 0
% local cospp3y = 0
%
% local function f(x,y)
%     if x == 1 then
%         cospp8y = cos(pp8*y)
%         cospp3y = cos(pp3*y)
%     end
%     local z = cos(pp8*x)*cospp3y + cos(pp3*x)*cospp8y
%     if z > 0 then
%         return '1'
%     else
%         return '0'
%     end
% end
%
% potrace.setbitmap("mybitmap", potrace.contourplot(N,N,f))
% \stopluacode
%
% \startMPpage[offset=1TS]
%     path p ; p := lmt_potraced [
%         stringname = "mybitmap",
%         value      = "1",
%         threshold  = 0.25,
%         optimize   = true,
%     ] ;
%
% fill p withcolor "darkred" ;
% \stopMPpage

% \startluacode
% local sin, cos, pi = math.sin, math.cos, math.pi
%
% local N   = 250
% local xx  = 6/N
% local yy  = 6/N
%
% local function f(x,y)
%     local x = xx*x - 3
%     local y = yy*y - 3
%     -- local z = y*x^2+2*y^2-3
%     -- local z = (x+y+1)^2 - 0.1
%     -- local z = x^2-y^2-1
%     -- local z = x^y - y^x
%      local z = -y*(1/(x^2+y^2)^(3/2)-1)
%     -- local z = 2*x*(x^2+y^2) - (3*x^2-y^2)
%     -- local z = x^2+y^2 - (x^2+y^2-1*x)^2 + y^3
%     -- local z = x*y^6 + y^7 + 2*x*y^3*sin(x) + 2*y^4*sin(x) + x*sin(x)*sin(x) + y*sin(x)*sin(x)
%     if z > 0 then
%         return '1'
%  -- elseif z < 0 then
%  --     return '2'
%     else
%         return '0'
%     end
% end
%
% potrace.setbitmap("mybitmap", potrace.contourplot(N,N,f))
% \stopluacode
%
% \startMPpage[offset=1TS]
%     path p ; p := lmt_potraced [
%         stringname = "mybitmap",
%         value      = "1",
%         % explode    = true,
%         % tolerance  = 0.5,
%         threshold  = 1.5,
%         optimize   = true,
%     ] ;
%     p := p scaled 4 ;
%     fill p withcolor "darkred" withtransparency (1,0.5) ;
%     draw p withpen pencircle scaled 5 withcolor "darkred" ;
%     drawpoints p withcolor white ;
%     currentpicture := currentpicture xsized 12cm ;
% \stopMPpage
